\documentclass[10pt,a4paper]{amsart}
\usepackage[margin=19mm]{geometry}
\usepackage{amsmath,amssymb,mathtools}
\usepackage[scr=rsfs,cal=euler]{mathalfa}
\usepackage{xfrac}
\usepackage[unicode,hidelinks,bookmarks=false]{hyperref}
\newcommand{\ie}{i.e.\ }
\newcommand{\cf}{cf.\ }
\newcommand{\resp}{respectively\ }
\newcommand{\Subsection}{\subsection}
\providecommand{\nobreakdash}{\nobreak\hbox{-}}
\setlength{\emergencystretch}{2em}
\multlinegap0pt
\usepackage{bm}
\usepackage{bbm}
\usepackage{dsfont}
\usepackage{xspace}
\usepackage{cancel}
\usepackage{mathtools}
\usepackage{amsthm}
\makeatletter
\@ifundefined{theorem}{\newtheorem{theorem}{Theorem}}{}
\makeatother
\makeatletter
\expandafter\ifx\csname proof\endcsname\relax
\newenvironment{proof}{\prepf\rm}{\endprepf}
\fi
\makeatother
\title[A new look at twin reduction]{A new look at twin reduction}
\author{Peter J. Cameron}
\date{Source: arXiv:2507.10189v1 (CC BY 4.0)}
\begin{document}
\maketitle

\noindent\emph{Source and licence.} A new look at twin reduction. Version arXiv:2507.10189v1 (CC BY 4.0), distributed under the Creative Commons Attribution 4.0 International licence, as recorded in the arXiv licence metadata for this exact version and shown on the abstract page https://arxiv.org/abs/2507.10189v1. Full rights notice: licenses/arxiv-2507.10189-CC-BY-4.0.txt.\par\smallskip

\noindent\textit{Editorial scope.} Counted unit: the Theorem whose source label is \texttt{None} in the article (source0.tex lines 148--155, source label \texttt{None}), delivered together with the article's own complete original proof of it (source0.tex lines 157--216, one \texttt{proof} environment of 60 lines). No mathematics is written, repaired or re-derived: statement and proof are the article's own text line for line. Required context retained uncounted: none. Every definition, display and label of those lines is retained unchanged. Omitted from this package and not redistributed: 1--147, 156--156, 217--308. Nothing inside a retained excerpt is omitted or altered: every retained body is delivered exactly as the article's lines are written, so no omission receipt of any kind accompanies this package. Citations: the retained excerpts carry no citation command at all, so no bibliography entry of the article is transcribed and no reference is redistributed. Reference adaptation: none was needed and none was made; every reference inside the delivered unit resolves inside it, so no unresolved reference or citation is produced. Printed slips kept literal: no printed word, symbol, hypothesis or number of the counted statement or of its proof was altered. Layout and class: the article's own class is \texttt{article} with packages including ; the wrapper is the lane's pinned \texttt{amsart} editorial wrapper at 10pt on A4, which restates the theorem-like environments the retained text needs and adds the article's own macro definitions that the retained text needs (none), with extra wrapper packages (bm, bbm, dsfont, xspace, cancel, mathtools). The delivered numbering is the unit's own. Assets: no retained span contains a figure environment, a graphics-inclusion command, a PDF-page inclusion, a table or a TikZ picture; no image file is copied into this package, the package declares no asset dependency and no image file is redistributed with it. Licence: version arXiv:2507.10189v1 of this article is distributed under the Creative Commons Attribution 4.0 International licence, as recorded in the arXiv licence metadata for this exact version and shown on the abstract page https://arxiv.org/abs/2507.10189v1. A full rights notice is delivered at licenses/arxiv-2507.10189-CC-BY-4.0.txt. Characters: every retained excerpt is pure ASCII, so the delivered engine, XeLaTeX with its default Latin Modern text font, reports no missing character.
\begin{theorem}
Let $\Gamma$ be a graph.
\begin{itemize}
\item[(i)] The join of any two sibling partitions is a sibling partition.
\item[(ii)] The unique maximal sibling partition is the partition produced by
complete twin reduction on $\Gamma$.
\end{itemize}
\end{theorem}

\begin{proof} (i) Let $\Pi_1$ and $\Pi_2$ be sibling partitions, and let
$\Pi=\Pi_1\vee\Pi_2$. We show that $\Pi$ satisfies (a) and (b). We begin
with (b). Let $x$ and $y$ lie in the part $A$ of $\Pi$. By definition, there
is a sequence $x=x_1,x_2,\ldots,x_m=y$ of vertices in $A$ such that each
consecutive pair is contained in a part of either $\Pi_1$ or $\Pi_2$. Now
take a vertex $z\notin A$. Since $\Pi_1$ and $\Pi_2$ satisfy (b), if
$z\sim x=x_1$, then $z\sim x_2$, \dots, and $z\sim x_m=y$. The argument is
similar if $z\not\sim x$. So (b) holds.

\smallskip

Now we prove (a). Suppose that it is false. Then there is a $4$-vertex path,
say $(a,b,c,d)$, contained in $A$. First we prove the following

\subparagraph{Claim} The points $a,b,c,d$ lie in different parts of $\Pi_1$ and
in different parts of $\Pi_2$.

Since $\Pi_1$ and $\Pi_2$ satisfy (a), the four points cannot all lie in the
same part of either.

Suppose that three of the points lie in the same part of $\Pi_1$ but the fourth
does not. Then, since $\Pi_1$ satisfies (b), the fourth point is joined to all
or none of the other three, which is clearly false. The same holds for $\Pi_2$.

Suppose that two of the points lie in the same part $A_1$ of $\Pi_1$ but the
other two do not lie in $A_1$. Then the two points not in $A_1$ are joined to
the same points in $A_1$, which (from the structure of the path) is not
possible. The same holds for $\Pi_2$. 

So the claim is proved.

\smallskip

Now $a$ and $d$ lie in the same part of $\Pi$, so there is a sequence
$a=x_1,x_2,\ldots,x_m=d$ where each consecutive pair lies in a part of
either $\Pi_1$ or $\Pi_2$. We can assume that the path of length $3$ is
chosen to minimise the length of such a sequence (which cannot be zero).
Now $a$ and $x_2$ lie in the same part of, say, $\Pi_1$ but $b,c,d$ do not
lie in this part; so $a$ and $x_2$ have the same adjacencies to $b,c,d$.
In particular, this means that $(x_2,b,c,d)$ is a path of length $3$, with
a shorter sequence joining its endpoints.

This contradiction shows that no path of length $3$ can exist in a part of
$\Pi$, so (b) is proved.

\medskip

Now we prove (ii). Let $\Pi$ be the maximal sibling partition (which is unique
by (i)). We claim that vertices
in different parts of $\Pi$ cannot be twins. For if they were, then replacing
these two parts by their union would give a coarser partition which is also a
sibling partition. (For the disjoint union and the sum of cographs are cographs;
and if $a\in A$ and $b\in B$ are twins, then they all vertices in $A\cup B$
have the same neighbours outside this set.)

So the quotient of $\Gamma$ by $\Pi$ is twin-reduced.

On the other hand, the remark before the theorem shows that we can reach $\Pi$
by twin reduction; but we cannot go any further.
\end{proof}

\end{document}
